% ==============================================================================
% KATA PENGANTAR (BERWARNA RESMI BPS)
% ==============================================================================

\begin{center}
  {\Large \textbf{\textcolor{bpsnavy}{Kata Pengantar}}}
\end{center}
\vspace*{0.6cm}

Puji dan syukur kami panjatkan kehadirat Tuhan Yang Maha Kuasa, atas rahmat dan karunia-Nya sehingga kami dapat menyusun dan menyelesaikan \textbf{Laporan Pembinaan Desa Cinta Statistik (Desa Cantik)} di Desa Popontolen, Kecamatan Tumpaan Kabupaten Minahasa Selatan dapat diselesaikan dengan tepat pada waktunya.

Laporan ini disusun untuk memenuhi rangkaian Pembinaan Desa Cantik 2026 yang diselenggarakan oleh Badan Pusat Statistik (BPS). Proses penyusunan laporan ini tidak terlepas dari peran banyak pihak yang telah membantu proses pembinaan Desa Cantik di Desa Popontolen. Oleh karena itu, pada kesempatan ini kami mengucapkan terima kasih yang sebesar-besarnya kepada seluruh pihak yang turut serta dalam pembinaan Desa Cantik ini.

Kami menyadari bahwa penyusunan laporan ini masih jauh dari kesempurnaan. Kritik dan saran yang membangun sangat kami harapkan untuk perbaikan tahap pembinaan Desa Cantik berikutnya. Demikian, semoga laporan ini dapat bermanfaat untuk kemajuan BPS.

\vspace*{1.5cm}
\begin{flushright}
  \begin{minipage}{0.58\textwidth}
    \centering
    Amurang Barat, 10 Agustus 2026\\[0.1cm]
    \textbf{\textcolor{bpsnavy}{Kepala BPS Kabupaten Minahasa Selatan}}\\[2.5cm]
    \textbf{\underline{\textcolor{bpsdark}{Irena Listianawati, SST, SE, M.Si}}}\\[0.1cm]
    \textcolor{gray!80}{\footnotesize Pembina Tingkat I / BPS Kab. Minahasa Selatan}
  \end{minipage}
\end{flushright}
